%%
%% MICRO 2026 Paper: Model-Defined Remapper
%% Based on the official MICRO 2026 LaTeX template (acmart sigconf)
%%
%TC:macro \cite [option:text,text]
%TC:macro \citep [option:text,text]
%TC:macro \citet [option:text,text]
%TC:envir table 0 1
%TC:envir table* 0 1
%TC:envir tabular [ignore] word
%TC:envir displaymath 0 word
%TC:envir math 0 word
%TC:envir comment 0 0

\documentclass[sigconf, screen, review]{acmart}

\AtBeginDocument{%
  \providecommand\BibTeX{{%
    Bib\TeX}}}

\setcopyright{none}
\copyrightyear{2026}
\acmYear{2026}
\acmDOI{XXXXXXX.XXXXXXX}
\acmConference[MICRO 2026]{The 59th IEEE/ACM International Symposium on Microarchitecture}{October 31--November 04, 2026}{Athens, Greece}
\acmISBN{978-X-XXXX-XXXX-X/XX/XX}

%%\acmSubmissionID{123-A56-BU3}

\settopmatter{printfolios=true}
\settopmatter{printacmref=false}

\begin{document}

%% ============================================================
%% TITLE
%% ============================================================
\title{Model-Defined Remapper: Virtual-to-Physical Topology Mapping for Chiplet-Based Architectures}
\subtitle{\normalsize{MICRO 2026 Submission
    \textbf{\#NaN} -- Confidential Draft -- Do NOT Distribute!!}}

%% ============================================================
%% ABSTRACT
%% ============================================================
\begin{abstract}

% TODO: Write abstract (~150-250 words)
% Structure: (1) Problem context, (2) Key challenge, (3) Our approach/insight, (4) Key results

\end{abstract}

\keywords{topology mapping, chiplet, network-on-chip, virtualization, reconfigurable architecture}

\maketitle

%% ============================================================
%% 1. INTRODUCTION (1.5-2 pages)
%% ============================================================
\section{Introduction}
\label{sec:intro}

% TODO: Paragraph 1 - Big picture motivation (chiplet era, heterogeneous integration)

% TODO: Paragraph 2 - The problem (virtual-to-physical topology mapping gap)

% TODO: Paragraph 3 - Why existing solutions fall short

% TODO: Paragraph 4 - Key insight / our approach

% TODO: Paragraph 5 - Contributions list
% Our key contributions are as follows:
% \begin{itemize}
%   \item We identify and formalize the virtual-to-physical topology mapping problem for chiplet-based architectures.
%   \item We propose \textit{Model-Defined Remapper}, a framework that ...
%   \item We evaluate ... across ... topologies and ... workloads, demonstrating ...
% \end{itemize}

%% ============================================================
%% 2. BACKGROUND & MOTIVATION (1-1.5 pages)
%% ============================================================
\section{Background and Motivation}
\label{sec:background}

\subsection{Chiplet-Based Architectures}
\label{sec:bg-chiplet}
% TODO: Background on chiplet integration, interposer networks, NoC topologies

\subsection{Virtual vs. Physical Topology}
\label{sec:bg-topology}
% TODO: Define virtual topology (application/model view) vs. physical topology (hardware)

\subsection{Motivating Example}
\label{sec:motivation}
% TODO: Concrete example showing performance gap between naive and optimized mapping
% Include a motivating figure showing the problem

%% ============================================================
%% 3. DESIGN / ARCHITECTURE (3-4 pages)
%% ============================================================
\section{Model-Defined Remapper}
\label{sec:design}

\subsection{Overview}
\label{sec:design-overview}
% TODO: High-level architecture diagram and design philosophy
% Figure: System overview of Model-Defined Remapper

\subsection{Problem Formulation}
\label{sec:formulation}
% TODO: Formal definition of the mapping problem
% Show NP-hardness or computational complexity argument

\subsection{Mapping Algorithm}
\label{sec:algorithm}
% TODO: Core algorithm / heuristic for topology mapping

\subsection{Hardware Support}
\label{sec:hardware}
% TODO: Any required hardware extensions or reconfiguration mechanisms

%% ============================================================
%% 4. IMPLEMENTATION (0.5-1 page)
%% ============================================================
\section{Implementation}
\label{sec:implementation}

% TODO: Simulation framework, tools used, RTL/cycle-accurate details
% Describe how the remapper is implemented/integrated

%% ============================================================
%% 5. EVALUATION (2.5-3 pages)
%% ============================================================
\section{Evaluation}
\label{sec:evaluation}

\subsection{Methodology}
\label{sec:methodology}
% TODO: Simulation setup, baselines, benchmarks/workloads, metrics
% Table: Simulation configuration parameters

\subsection{Performance Results}
\label{sec:results-perf}
% TODO: Latency, throughput comparisons vs. baselines
% Figures: Performance comparison charts

\subsection{Energy and Area Analysis}
\label{sec:results-energy}
% TODO: Energy consumption, area overhead analysis

\subsection{Scalability}
\label{sec:scalability}
% TODO: Results across different topology sizes and configurations

\subsection{Sensitivity Analysis}
\label{sec:sensitivity}
% TODO: How results change with key parameters

%% ============================================================
%% 6. RELATED WORK (0.5-1 page)
%% ============================================================
\section{Related Work}
\label{sec:related}

\subsection{NoC Topology Design}
% TODO: Prior work on NoC topology design and optimization

\subsection{Topology Mapping and Virtualization}
% TODO: Prior work on application-to-topology mapping, NPU virtualization

\subsection{Chiplet Integration}
% TODO: Prior work on chiplet architectures and interposer design

%% ============================================================
%% 7. CONCLUSION (0.5 page)
%% ============================================================
\section{Conclusion}
\label{sec:conclusion}

% TODO: Summary of contributions and key results
% TODO: Future directions

%% ============================================================
%% REFERENCES
%% ============================================================
\bibliographystyle{ACM-Reference-Format}
\bibliography{references}

\end{document}
